\subsection{Compatible joint sections and their full intrinsic central fibres}

\begin{statement}{E.85}{Proved}{The compatible joint family and an exact index-nine central comparison}
Retain the actual local source, conjugate maps and
tangent of E.77--E.80. Write
\[
 V=N^{\rm ab},\qquad A_2=\gamma_2N/\gamma_3N,
 \qquad N_2=N/\gamma_3N,\qquad
 A=\ker(\mathcal J_2\to\mathcal J_1)=\mathbb Z_3^2.
 \tag{E85.1}
\]
The two coordinates of $A$ are the exponents of
$(c,1)$ and $(1,c)$ in E.80. Its actual operators
are $49I$ and the factor swap. Use the unchanged
integral cubic basis
$(f_1,f_2,A_0,A_1,A_2,B_0,B_1,B_2,C_1,C_2)$.

There is an actual full intrinsic local section
with E.77 values $(u_0,d)$ such that
\[
 \begin{gathered}
 J(d)=d^{-1},\qquad u_0T(d)u_0^{-1}=d^7,\qquad
 (h_*,\bar h_*)(d)=(y,y),\\
 (h_*,\bar h_*)(u_0)\bmod\gamma_3F^2
       =(x^{5/2}c^2,\ x^{1/2}y^{-2}c^3).
 \end{gathered}
 \tag{E85.2}
\]
Here $d$ is the positive inertia of the rational
cusp $B_0$. The proof constructs a fixed path and
its actual normalizing conjugators; the constants
two and three are forced by the full cusp identity.
Let $\alpha_0$ be the abelian class of this section
in $H^1(\mathbb Q_3,V)$ and $\beta_0$ its image
in $H^1(\mathbb Q_3,\mathcal J_1)$.

For E.80's whole family
\[
 a=y,\qquad a'=y^{-1},\qquad
 u=x^{5/2}c^N,\qquad u'=x^{-1/2}y^2c^M,
                   \qquad N,M\in\mathbb Z_3,
 \tag{E85.3}
\]
the actual intrinsic class-two preimage is nonempty
exactly for
\[
                       N+M\equiv2\pmod3.
 \tag{E85.4}
\]
Every compatible member already has a preimage over
the single intrinsic abelian class $\alpha_0$.
The complete geometric-conjugacy classes of the
compatible target family have the coordinate and
inverse
\[
       t=N-M+1\in\mathbb Z_3,\qquad
       N=2+t/2,\quad M=3-t/2.
 \tag{E85.5}
\]
In particular this is an exact existence converse
to E.83 for this family.

The entire fibre over $\alpha_0$, including variations
of both relative source generators, is evaluated as
follows. Put
\[
 \begin{gathered}
 \mathcal P=A_2/(I-T_2)A_2^{J_2}
       \simeq\mathbb Z_3^{24}\oplus C_3^9\oplus C_9^5,\\
 \mathcal K_2=\ker((I-J_2)(T_2-7I):A_2\to A_2)
       \simeq\mathbb Z_3^{22},\\
 \mathfrak F_{\alpha_0}:=
  \bigl(H^1(\mathbb Q_3,N_2)\to H^1(\mathbb Q_3,V)\bigr)^{-1}(\alpha_0)
       \xrightarrow{\ \sim\ }\mathcal P\oplus\mathcal K_2.
 \end{gathered}
 \tag{E85.6}
\]
The inverse of the last arrow is the actual section
with values
\[
             (u_0w_z,\ d w_a),
                 \qquad ([z],a)\in\mathcal P\oplus\mathcal K_2,
 \tag{E85.7}
\]
where $w_z,w_a$ are the ordered Hall representatives
modulo $\gamma_3N$. All four source values and the
trivial value on the order-two source subgroup are
then determined by E.77. The residual gauge is exactly
$z\mapsto z+(I-T_2)b$, $b\in A_2^{J_2}$; it leaves
$a$ unchanged.

The target fibre over $\beta_0$ is
\[
 \begin{gathered}
 \mathfrak G_{\beta_0}:=
  \bigl(H^1(\mathbb Q_3,\mathcal J_2)
          \to H^1(\mathbb Q_3,\mathcal J_1)\bigr)^{-1}(\beta_0)
       \xrightarrow{\ \sim\ }
             \mathbb Z_3\oplus\mathbb Z_3\oplus C_3,\\
   (\rho,\sigma,\epsilon)
        =(n-m,\ s,\ (n+m)/2\bmod3).
 \end{gathered}
 \tag{E85.8}
\]
In the last line the central changes from the
origin (E85.2) are $(n,m)$ on its horizontal value
and $(s,s)$ on its relative value. For
$\epsilon\in\{0,1,2\}$ an actual inverse is
\[
 \begin{gathered}
 u_0^{\mathcal J}
       (c^{\rho/2+\epsilon},c^{-\rho/2+\epsilon}),\qquad
 a_0^{\mathcal J}(c^\sigma,c^\sigma),\qquad
 a_0^{\mathcal J}=(y,y).
 \end{gathered}
 \tag{E85.9}
\]

Let $F_2:A_2\to\mathbb Z_3^2$ be the two-by-forty-five
matrix obtained from the four rows of (E83.11): on
a Hall vector $[q_a,q_b]$, its two entries are
\[
 Q_{\epsilon,a}Q_{\epsilon+1,b}
       -Q_{\epsilon+1,a}Q_{\epsilon,b},\qquad \epsilon=0,2.
 \tag{E85.10}
\]
Its exact integral image and arithmetic comparisons are
\[
 \begin{gathered}
 L_2=\{(n,m):n+m\in3\mathbb Z_3\},\qquad
 F_2T_2=49F_2,\quad F_2J_2=\operatorname{swap}F_2,\\
 F_2(\mathcal K_2)=\{(3s,3s):s\in\mathbb Z_3\}.
 \end{gathered}
 \tag{E85.11}
\]
Consequently the complete natural map on the two
fibres in (E85.6) and (E85.8) is
\[
 \begin{gathered}
 ([z],a)\longmapsto
       \bigl((F_{2,1}-F_{2,2})z,\ F_{2,1}a,\ 0\bigr),\\
 \operatorname{im}(\mathfrak F_{\alpha_0}\to\mathfrak G_{\beta_0})
       =\mathbb Z_3\oplus3\mathbb Z_3\oplus0,\\
 \operatorname{coker}
       \bigl(H^1(\mathbb Q_3,A_2)\to H^1(\mathbb Q_3,A)\bigr)
       \simeq C_3^2,\\
 \ker\bigl(H^1(\mathbb Q_3,A_2)\to H^1(\mathbb Q_3,A)\bigr)
       \simeq\mathbb Z_3^{44}\oplus C_3^9\oplus C_9^5.
 \end{gathered}
 \tag{E85.12}
\]
Every point in the displayed image has exactly this
compact fibre, with the explicit coordinates below.
Thus each compatible member (E85.3) has a complete
preimage \emph{over $\alpha_0$} of this type.

There are small literal splitting vectors. With
brackets in the retained Hall convention put
\[
 \begin{gathered}
 z_*=-\tfrac12[f_2,f_1],\qquad
 a_*=-[B_0,f_1]+[B_0,f_2]-[B_0,A_1],\\
 F_2z_*=(1/2,-1/2),\qquad
 a_*\in\mathcal K_2,\qquad F_2a_*=(3,3).
 \end{gathered}
 \tag{E85.13}
\]
The inverse to the image coordinates $(\rho,3s,0)$
can therefore be taken to be the section (E85.7)
with $z=\rho z_*$, $a=s a_*$. Adding the exact
kernel coordinates gives all its preimages over
$\alpha_0$.

For completeness these kernel coordinates use the
actual retained Smith and relative-kernel maps.
Let $L_GGR_G=D_G$ be E.82's exact Smith system for
$G=(I-T_2)(I+J_2)$, where
\[
 D_G=\operatorname{diag}(1^7,3^9,9^5,0^{24}).
 \tag{E85.14}
\]
Write $p=L_Gz$. Set its first seven entries to zero,
take $p_8,\ldots,p_{16}$ in $C_3$ and
$p_{17},\ldots,p_{21}$ in $C_9$, and retain its final
twenty-four entries in $\mathbb Z_3$. In that order
the row for $\rho$ on the final twenty-four entries is
\[
 r=(-1,-1,-1,0,0,0,0,-3,3,0,-1,0,-1,0,0,3,-3,0,0,-2,-3,-3,-1,0).
 \tag{E85.15}
\]
Let $B$ and $B^{\ell}$ be E.75's retained saturated
kernel basis and left inverse for $\mathcal K_2$;
write $a=B\eta$, $\eta=B^{\ell}a\in\mathbb Z_3^{22}$.
The exact row for $\sigma/3$ is
\[
 s=(0,0,0,0,-3,0,0,0,0,0,1,0,-3,2,-2,0,0,0,-1,2,-1,-2).
 \tag{E85.16}
\]
For any image value $(\rho,3v,0)$ retain all the
fourteen cyclic coordinates, the twenty-three
entries $p_{23},\ldots,p_{45}$, and the twenty-one
entries $\eta_j$ with $j\ne11$. The inverse inserts
\[
 \begin{gathered}
 p_{22}=\sum_{j=2}^{24}r_j p_{21+j}-\rho,\qquad
 \eta_{11}=v-\sum_{j\ne11}s_j\eta_j,\\
             z=L_G^{-1}p,\qquad a=B\eta,
 \end{gathered}
 \tag{E85.17}
\]
and then applies (E85.7). These are mutually inverse
homeomorphisms on the full geometric-conjugacy
fibres. The source generators are E.77's actual
generators; $B$ is used as a basis of its displayed
coefficient kernel, without identifying distinct
earlier choices of relative source elements.

Restricting to $a=0$ in (E85.7) recovers the
embedded forty-five-parameter cusp family. Its
map onto the compatible target family has exact
fibres
\[
                   \mathbb Z_3^{23}\oplus C_3^9\oplus C_9^5.
 \tag{E85.18}
\]
The entire intrinsic preimage can also contain
sections over other abelian preimages of $\beta_0$.
The source of the complete fibre assertion in
(E85.12) is explicitly $\mathfrak F_{\alpha_0}$;
no other intrinsic abelian class is identified
with $\alpha_0$.

\emph{Manuscript: \S51.16, folios733v--736v,
scans172--178, formulas(371)--(379), Part E(c),(d)
and follow-up(4). This evaluates actual compatible
joint preimages and their full central fibres over
a specified intrinsic abelian section.}
\end{statement}

\begin{proof}
\emph{The actual rational cusp origin.}
For the path construction use actual global complex conjugation
$c_\infty\in G_{\mathbb Q}$, which fixes the two rational unit
 tangents. Its image in the global quotient $\Gamma$ is the marked
$c$, and its based action on $N$ is exactly $J$ by E.34.3 and
E.78's rational tangent transport. The local involution of E.77
has the same image under $\mathcal D\to\Gamma$, as constructed
in E.74. Choose an actual global lift $\widetilde\tau$ of the
marked $\tau\in P$; its based action is $T$ and its three-adic
cyclotomic value is seven. Thus the global path actions used below
induce the prescribed based operators. The element
$p\widetilde\tau(p)^{-1}$ supplies the displayed norm identity,
and E.77 constructs the local section from that identity and
$J(d)=d^{-1}$. This is the exact passage from the global path
groupoid to the local section presentation.

Choose a rational unit tangent at $B_0$ fixed by
the order-two action, and start with the torsor
of geometric pro-three paths from the stated
base tangent to it. This torsor has a fixed
point. Here the needed coprime fixed-point fact
has a direct proof: for a finite three-group,
every order-two cocycle is a coboundary. Induct
on the group through an invariant elementary
abelian central subgroup; on that subgroup
averaging divides by two. For a pro-three group,
the nonempty solution sets in finite invariant
quotients give a nonempty compact inverse limit.
The same argument proves that an equivariant
surjection of pro-three groups is onto on fixed
points: the defect of a lift is such a cocycle
in its kernel.

Use a fixed path $p$, and put
$d=p\iota_{B_0}p^{-1}$. Thus $J(d)=d^{-1}$.
The maps $h,\bar h$ have order one at the
respective target cusp $1$, so the first image
is $\alpha y\alpha^{-1}$ for an actual
$\alpha\in F$. The equality $J(d)=d^{-1}$ and
the literal factor-swap involution of E.80
then make its second image exactly
$I(\alpha)yI(\alpha)^{-1}$. Set
\[
       w=y^{-\ell(\alpha)}\alpha^{-1}.
 \tag{E85.19}
\]
Then $\ell(w)=0$. Hence $(w,Iw)$ belongs to
$\mathcal J^{\jmath}$ and conjugates this inertia
pair to $(y,y)$. The actual equivariant
surjection $N\to\mathcal J$ in E.80 lifts it
to a fixed element of $N$, by the preceding
fixed-point argument. Change $p$ by that
element. This retains a fixed path and proves
the exact inertia assertion in (E85.2).

Initially put $u=pT(p)^{-1}$. The actual tangent
action gives $uT(d)u^{-1}=d^7$ in the full
group. Project this identity to each class-three
target, retaining all five Hall coordinates.
The class-two coefficient first forces the two
$x$-coordinates of $u$ to be $5/2,1/2$.
Membership of its pair in $\mathcal J$ makes
their $y$-coordinates $V,V-2$ for one
$V\in\mathbb Z_3$. The actual E.78 polynomial
identity, with central coordinates $r_i,s_i$
retained, is
\[
 \begin{aligned}
 T_{-5/2}:\quad
  x^{5/2}y^Vc^{\,2-5V/2}d^{r_1}e^{s_1}
       &\quad\hbox{conjugates }T_{-5/2}(y)\hbox{ to }y^7,\\
 T_{-1/2}:\quad
  x^{1/2}y^{V-2}c^{\,3-V/2}d^{r_2}e^{s_2}
       &\quad\hbox{conjugates }T_{-1/2}(y)\hbox{ to }y^7.
 \end{aligned}
 \tag{E85.20}
\]
Conversely the two $c$-coordinates displayed
are forced by those identities: their errors
are respectively
$e^{7(2-N-5V/2)}$ and
$e^{7(3-M-V/2)}$.
Replace $u$ in the full group by $u_0=d^{-V}u$.
This preserves the norm identity and gives
exactly the two class-two words in (E85.2).
No value of $r_i,s_i$ is assigned or needed.
E.77 now supplies the actual full intrinsic
section with values $u_0,J(u_0),d,J(d)$.

\emph{Every compatible family member.}
Let $e=[f_2,f_1]$ and $f=[B_0,f_1]$ in $A_2$.
Their actual images, computed from (E83.11), are
\[
                   F_2e=(-1,1),\qquad F_2f=(-2,-1).
 \tag{E85.21}
\]
They form an integral basis for $L_2$. For
$n=N-2$, $m=M-3$ and $\ell=(n+m)/3$, the
literal Hall vector
\[
                        z=(2\ell-n)e-\ell f
 \tag{E85.22}
\]
has image $(n,m)$ whenever (E85.4) holds.
Since $A_2$ is central in $N_2$, the pair
$(u_0w_z,d)$ still has norm defect one. Its
joint image is precisely (E85.3), with the
conjugate source values prescribed in E.80.
This proves existence. E.83 proves emptiness
for all incompatible members, completing the
equivalence (E85.4).

For the complete target gauge, a normalized
geometric conjugator is an element of
$\mathcal J_2^{\jmath}$. Comparing the fixed
abelian horizontal coordinates of two members
forces its abelian vector to be zero, since
$7-1$ is injective on that lattice. Therefore
it is $(c^b,c^b)$, $b\in\mathbb Z_3$.
Its exact change is
\[
                       (N,M)\longmapsto(N-48b,M-48b).
 \tag{E85.23}
\]
For compatible members the actual normalizing
value $b=(N+M-5)/96$ is integral. It gives
(E85.5), and no nonzero remaining $b$ fixes
the resulting pair. Conversely equality of
$t$ for compatible members makes their common
coordinate difference divisible by three,
so (E85.23) supplies the required integral
conjugator. This proves the complete orbit
statement, including the nonunit factor in
the normalizing denominator.

\emph{The full fibre over the chosen intrinsic abelian class.}
The abelian section $\alpha_0$ has no geometric
stabilizer. Indeed the actual $T_1$ on $V$ has
eight eigenvalues seven and an elliptic block
with characteristic polynomial $X^2+X+7$;
neither has eigenvalue one. Thus its fixed
lattice is zero. An intrinsic section with
abelian class $\alpha_0$ can consequently be
aligned with its actual abelian representative
by a unique abelian conjugator. After normalizing
the order-two source value to one, that
conjugator is $J_1$-fixed. Fixed-point lifting
through $N_2\to V$ gives an actual fixed lift.
All remaining conjugators lie in $A_2^{J_2}$.

The aligned values therefore have the unique
form (E85.7). Centrality and the full identity
in (E85.2) give its literal norm value
\[
 (u_0w_z)T(dw_a)(u_0w_z)^{-1}(dw_a)^{-7}
                            =(T_2-7I)a\quad\hbox{in }A_2.
 \tag{E85.24}
\]
E.77 makes the source relation precisely
$(I-J_2)(T_2-7I)a=0$. Conjugation by a remaining
$b\in A_2^{J_2}$ changes $z$ by $(I-T_2)b$ and
does not change $a$. This proves every direction
of (E85.6)--(E85.7). The Smith and saturated
kernel maps used are exactly those of E.82 and
E.75 for these actual coefficient operators.
The E.78 Frobenius transport preserves them by
the explicit CM and gluing comparison in E.83.

Apply the same argument to the target abelian
lattice. Its fixed group is zero because its
operator is $7I$. For its central lattice $A$
the condition in (E85.24) is
$42(I-\operatorname{swap})(s_1,s_2)=0$.
It is equivalent to $s_1=s_2=s$. The remaining
central conjugator is $(b,b)$, changing
$(n,m)$ by $(-48b,-48b)$. Thus its coordinates
are exactly (E85.8). To check the inverse,
take $\epsilon$ to be the residue of
$(n+m)/2$ and use
\[
               b=\frac{(n+m)/2-\epsilon}{48}\in\mathbb Z_3.
 \tag{E85.25}
\]
This is the actual conjugator producing
(E85.9), and proves that these coordinates
already include full geometric conjugacy.

\emph{The coefficient map and all inverse fibre coordinates.}
The coefficient map of the central extensions
is the alternating map (E85.10), since both
targets use $c=[x,y]$. Its two-by-two image
minor (E85.21) has determinant three. Every
column has sum divisible by three, proving
the exact image $L_2$. The two arithmetic
intertwiners in (E85.11) follow from the actual
four unit characters and are also checked on
all forty-five Hall columns.

For $a\in\mathcal K_2$ those intertwiners give
\[
       42(I-\operatorname{swap})F_2a=0.
 \tag{E85.26}
\]
Thus $F_2a$ is diagonal. Its membership in
$L_2$ makes its common entry divisible by
three. The retained vector $a_*$ in
(E85.13) satisfies the defining integral
kernel equation and has image $(3,3)$.
This proves equality in (E85.11), not merely
a finite-level bound. Similarly the vector
$z_*$ has the displayed image and gives all
values of $\rho$. The third target coordinate
vanishes because every $F_2z$ has coordinate
sum divisible by three. These facts prove
the exact map and image (E85.12) with the
explicit section (E85.13).

For the complete coordinate inverse use the
literal matrices in (E85.14). Multiplying
$(F_{2,1}-F_{2,2})L_G^{-1}$ gives zero on
the first twenty-one columns and the row
(E85.15) on the last twenty-four. Thus the
map kills all fourteen torsion coordinates,
and its first free coefficient is minus one.
The exact product $F_2B$ has two identical
rows, equal to three times (E85.16); its
eleventh coefficient after division by three
is one. Solving these two unit equations
gives precisely (E85.17). The retained
matrices include both inverse pairs over
$\mathbb Z_{(3)}$, and $B^{\ell}B=I$.
Their ranks and these unit coefficients
therefore give the asserted kernel type,
with no unproved index or saturation step.

Restricting the relative central coordinate
$a$ to zero removes its twenty-two free
coordinates. The remaining kernel of the
surjective row $\rho:\mathcal P\to\mathbb Z_3$
is (E85.18), with inverse obtained by the
same first equation in (E85.17). By contrast
the full fibre over $\alpha_0$ retains the
twenty-one free relative coordinates in
that formula. This is the exact inclusion
between the two source fibres.

Every space used here is compact, the
displayed groups are Hausdorff, and all
Hall representatives, matrix maps and
unit divisions are continuous. The explicit
inverse formulas prove the homeomorphism
assertions directly. In particular no
closure is silently added to a gauge image;
these images are already compact images
of the indicated finite free lattices.

The executed certificate independently
expands all ninety alternating tensor
images, checks the exact image lattice,
both arithmetic intertwiners, all inverse
Smith and kernel maps, and the two literal
splitting vectors. It verifies four symbolic
class-three cusp norm identities while
retaining the unknown central path entries.
It also checks all $729$ sampled target
family pairs, $1,296$ exact coordinate
inverse identities, and the index-nine
finite coordinate quotients through
modulus $81$. The complete input hashes,
matrices and actual output are retained.
\end{proof}
